\vspace{-0.5em}
\section{Conclusion}\label{sec:conclusion}
This paper presented \emph{PaTAS}, a Subjective Logic framework propagating trust in neural networks through a parallel structure of \emph{Trust Nodes} and \emph{Trust Functions} running alongside feedforward and backpropagation; its \emph{Parameter Trust Update} and \emph{Inference-Path Trust Assessment} capture how input quality, learned parameters, and activation paths jointly shape the trustworthiness of a prediction.

Evaluation on real-world and adversarial datasets shows that PaTAS detects trust degradation caused by data poisoning and adversarial patching while remaining interpretable and stable, providing reliability information that accuracy alone does not reflect, in particular when training-data provenance is unknown. In a comparison spanning output-only uncertainty baselines and dedicated input- and feature-space detectors, the detection of distribution violations is carried by the conformity-derived input opinions, which alone outperform every dedicated detector on the violation tasks (OOD $1.000$, corruption $1.000$ on MNIST, where feature-space methods fall to $0.59$--$0.74$); trust propagation contributes what no detector expresses: a test-data-free training audit, per-class trust gaps on poisoned models that persist without oracle patch knowledge, graceful composition where a single signal fails, and in-distribution selective prediction matching the softmax filter after identical calibration. Applicability to a dataset is predictable a priori from the conformity model's own statistics (\cref{tab:diagnostics}).

Several directions remain open. The evaluation used shallow fully-connected architectures on small benchmarks; extending trust propagation to deeper convolutional and transformer models, and enriching the input-trust construction beyond marginal conformity so that unregistered imagery (GTSRB, CIFAR-10) moves inside the applicability boundary, are the primary next steps. Further studies would strengthen the framework: a sensitivity analysis under noisy or incorrect initial input and label opinions; a deeper study of the gradient-to-evidence mapping beyond the $\epsilon$ sensitivity of \cref{tab:cancer_results}; and an empirical measurement of runtime and memory overhead against the analytical $3\times$ estimate. The security scope could broaden from detection to provable robustness, with guarantees under bounded poisoning and evaluation against evasion, model-stealing, and membership-inference attacks. On the modeling side, the effect of activation functions on trust propagation is uncharacterized, and the gradient-counting used for \(T_{\theta_i \mid y_\text{batch}}\) fixes the uncertainty component by construction, motivating adaptive quantification schemes that also account for parameter magnitude.
